\section{The experimental problem}\label{sec:problem}
A preparation and a sequence of causal interventions determine which predictions have actually been acquired. Two different obstructions then arise. A formally reachable set may carry little acquired mass, and a good terminal codebook may have no equally good recursive implementation. The first obstruction concerns probability; the second concerns the loss of access to discarded reports. A foundations theorem must connect them without replacing either by a dimension count.

The central result below is Theorem~\ref{thm:inward}. It starts with a report kernel and a predictive update, constructs the finite-label machine from a static covering, and controls that machine under the \emph{true} report law. This last qualification matters: a drift inequality under the approximate filter's own report distribution does not control an approximation driven by physical data. The construction combines inward quantization with two raw moment inequalities. It does not assume an error recurrence for an already constructed encoder. It also avoids comparison of all intermediate occupation masses with a terminal reference law.

The geometry is supplied by the actual terminal score law. Theorem~\ref{thm:inward} gives an upper bound on a general metric state domain. Proposition~\ref{prop:submass} turns a small-ball estimate for only a positive acquired submeasure into a matching lower bound, allowing the rest of the acquired law to be singular. Proposition~\ref{prop:oracle} treats the opposite difficulty: a finite acquisition protocol can make the entire predictive law atomic even though the unresolved physical quantity has a continuous or singular law. Its exact decomposition prevents substituting an ideal law for the actual acquisition.

Two verifications carry the general theorem beyond uniform one-step contraction. Theorem~\ref{thm:filter} treats a hidden-state experiment with observed intermittent mixing. Informative identity-transition observations are isometries in log odds and can occur in arbitrarily long runs. A finite inward compander on the unbounded log-odds domain nevertheless attains the optimal label scale without a logarithmic loss. Theorem~\ref{thm:singular} treats finite acquisition in a nonhomogeneous Cantor experiment. Both mode changes transport an old unrevealed tail, deletion expands errors, and mode probabilities may vary with time or vanish. No invariant preparation comparison enters the proof. The finite acquired law is computed explicitly.

The two results instantiate the same kernel theorem; neither is a premise of it. The complementary allocation-weighted, finite-chart, reset-preserved, and continuation-information results are retained in the accompanying mathematical supplement with their original hypotheses. The present route is not a claim that one stability condition subsumes every experiment. It supplies a different constructive mechanism: direct raw-kernel control of a recursively retained state, with acquired geometry supplying the converse.

\subsection{Raw experiments and legal information}
An experiment consists of standard Borel spaces $X_t,A_t,Y_{t+1}$, a probability preparation $\mu$ on $X_0$, and jointly measurable nonnegative normalized kernels
\begin{equation}\label{eq:raw}
 K_t(dx',dy,dc\mid x,a).
\end{equation}
Here $c$ records the declared nonnegative resource increments and elapsed physical time. Positivity means nonnegativity and normalization, not a uniform lower bound on report probabilities. Dirac kernels and failed acquisitions are allowed. A legal strategy is a kernel $\beta_t(da\mid h_t)$ using only the visible history, known calibration, and independent randomization. The history includes actions, reports, failures, and every published cost. A clock or remaining budget that changes legal continuations is part of the visible interface.

The Bayesian assertions fix the preparation and a legal exploration strategy independently of the encoder. An unknown parameter can be included in the hidden preparation with a specified prior; this does not turn a prior-specific quotient into a parameter-uniform sufficient statistic. The comparison theorem in Section~\ref{sec:resources} has its own uniform-in-parameter quantifiers. We work at a fixed deterministic horizon. Independent exact holds are allowed as specified below. Stopping rules are legal objects of the experiment, but the main risk theorem is not silently extended to arbitrary stopping times.

The successive kernel integrals define the actual path law. Write $\cH_t=\sigma(h_t)$ for its visible filtration. An online encoder retains a label in an alphabet of cardinality at most $M$, updates it using the current report and declared current action, and cannot reread earlier reports. A checkpoint encoder may inspect $h_T$ once and then retain at most $M$ labels. Independent randomization is allowed in both classes. Persistent labels do not include a free history tape, mode register, or real-valued posterior register. Required controller and internal clock states are counted separately in a serial implementation and are not silently removed from its total state.

\subsection{Executable tests and a realized quotient}
An executable future test is a legal finite continuation followed by a bounded measurable function of its actual reports. A hidden-state function is a test only when the experiment includes a measurement producing the required readout. Two histories are predictively equivalent when they have the same continuation interface and every such conditional test expectation agrees. A countable determining family $f_1,f_2,\ldots$ yields the evaluation map
\begin{equation}\label{eq:eval}
 h\longmapsto \big(\E[f_n\mid h]\big)_{n\ge1}.
\end{equation}
A realization requires measurability of its image and a valid update, not just a formal equivalence relation. The following sufficient certificate makes the relevant version issue explicit.

\begin{proposition}[Compact continuous-instrument realization]\label{prop:quotient}
Let $\mathcal B$ be a compact metrizable domain of beliefs, and let $\mathcal V\subset C(\mathcal B;\R)$ be the separable closed linear span of the continuous executable-test evaluations and declared continuation-interface coordinates, including constants. Choose a countable uniformly dense family in $\mathcal V$ for the evaluation map. Suppose, for each declared action $a$, there is a full-support reference measure $m_a$ on a Polish report domain and a continuous nonnegative density $g_a(\pi,y)$. On the declared domain $D_a=\{(\pi,y):g_a(\pi,y)>0\}$ let the normalized belief update $B_a(\pi,y)$ be continuous. Assume that, for every $f\in\mathcal V$ and every declared $y$, the function
\[
 \pi\longmapsto g_a(\pi,y) f(B_a(\pi,y))
\]
extends continuously by its instrument value to $\mathcal B$ and belongs to $\mathcal V$. Suppose these continuous evaluations determine all executable tests, including their allowed continuations. Then the image $S$ of the evaluation map on $\mathcal B$ is compact metrizable, predictive equivalence is equality in $S$, and the report law and update descend to $S$. The update is pointwise well defined on the positive-density declared domain. If a measurable statistic $r:H\to Z$, with $Z$ standard Borel, measurably determines every evaluation, the realized quotient is a measurable function of $r$ on its image (and after any fixed measurable extension off that image).
\end{proposition}
\begin{proof}
Normalize the countable determining evaluations to be bounded by one. Their product map $E:\mathcal B\to[-1,1]^{\N}$ is continuous, so $S=E(\mathcal B)$ is compact metrizable and Borel. Its fibers agree on the dense determining set, hence on $\mathcal V$ by uniform approximation, and hence on the declared executable tests. Taking $f=1$ in the instrument closure shows that equal fibers have equal $g_a(\pi,y)$ for every declared report. For positive common density, equality of the density-weighted evaluations followed by division by that \emph{positive} number shows that their updated beliefs have equal $E$-images. Thus both the report kernel and normalized update descend. Continuity of a map constant on fibers of a compact-to-Hausdorff quotient supplies continuity of the descended instrument evaluations. On each positive-density neighborhood the same quotient argument applies to the update. Nothing is asserted by normalization at density zero; an extra prescribed null-report version would require an extra closure certificate.

For the last assertion, for each $n$ let $u_n:Z\to\R$ be a measurable factor for the $n$th evaluation. The product $u=(u_n)_n$ is measurable. Since $S$ is Borel, replace $u(z)$ by a fixed $s_0\in S$ whenever $u(z)\notin S$. This is a measurable $S$-valued extension and agrees with $E$ on all histories. Under completed, almost-sure conventions the countably many equalities are first imposed on one common full-measure history set. No factorization through an unspecified arbitrary measurable codomain is used.
\end{proof}

The raw realizations below also verify their quotients directly: a finite-dimensional posterior, or a mode and cylinder posterior, determines every legal continuation and is separated by a terminal test. The compact certificate is sufficient, not necessary; the stable metric used for quantization need not induce the compact quotient topology. In particular, log odds has an unbounded metric even though a finite-state posterior lies in a compact simplex.

\subsection{A common squared task}
Fix a scalar terminal report $Y\in[0,1]$, or a finite collection of separately executable scalar probes with prescribed positive weights. Expectations of the weighted sum of their squared losses define the vector task; no joint counterfactual observation is postulated. Use the corresponding Euclidean norm. Set
\begin{equation}\label{eq:score}
 P_T=\E[Y\mid\cH_T],\qquad \nu_T=\law(P_T),\qquad
 B_T=\E\norm{Y-P_T}^2.
\end{equation}
All statements below are componentwise for this interpretation. For a probability law $\nu$ on the task box define
\begin{equation}\label{eq:q}
 q_M(\nu)=\inf_{\#C\le M}\int \operatorname{dist}(p,C)^2\,\nu(dp).
\end{equation}
The infimum permits arbitrary centers; projection to the task box is harmless. The symbols $R_\cp(T,M)$ and $R_\on(T,M)$ denote the infimal risks in the two legal classes. The law $\nu_T$ always comes from the actual preparation and exploration, not an invariant reference chosen for analytic convenience.
